\documentclass{article}
\usepackage{amsmath,amssymb,amsthm}
\usepackage{algorithm}
\usepackage{algpseudocode}
\usepackage{bm}
\usepackage{hyperref}
\newtheorem{assumption}{Assumption}
\newtheorem{theorem}{Theorem}
\newtheorem{lemma}{Lemma}
\newcommand{\topk}[1]{\mathcal{S}_K(#1)}
\begin{document}

\section*{Top‑K Stochastic Gradient Descent (Top‑K SGD)}

\section*{Mathematical Formulation}
Let $f:\mathbb{R}^d\rightarrow\mathbb{R}$ be a (possibly non‑convex) loss function.  
At iteration $t$ a stochastic gradient $g_t$ is sampled,
\[
g_t \;=\; \nabla f(w_t;\,\xi_t),\qquad \xi_t\sim\mathcal{D},
\]
where $\mathcal{D}$ is the data distribution.  

Define the \emph{top‑$K$ sparsification operator}
\[
\topk{g}\;\in\;\mathbb{R}^d,\qquad
[\topk{g}]_i=\begin{cases}
g_i, & i\in\mathcal{I}_K(g),\\[2mm]
0,   & \text{otherwise},
\end{cases}
\]
where $\mathcal{I}_K(g)$ contains the indices of the $K$ components of $g$ with largest absolute value, i.e.
\[
\mathcal{I}_K(g)=\arg\!\max_{|S|=K}\sum_{i\in S}|g_i|.
\]

The Top‑K SGD update is
\begin{equation}
\boxed{\,w_{t+1}=w_t-\eta_t\,\topk{g_t}\,} \tag{1}
\end{equation}
with learning‑rate schedule $\{\eta_t\}_{t\ge0}$.

When the algorithm is run on $M$ parallel workers, each worker $m$ computes a local stochastic gradient $g_t^{(m)}$, applies $\topk{\cdot}$, and the server aggregates the sparse vectors by averaging:
\[
w_{t+1}=w_t-\frac{\eta_t}{M}\sum_{m=1}^{M}\topk{g_t^{(m)}} .
\]

\section*{Pseudocode}
\begin{algorithm}[H]
\caption{Top‑K SGD}
\begin{algorithmic}[1]
\State \textbf{Input:} initial point $w_0\in\mathbb{R}^d$, stepsizes $\{\eta_t\}$, sparsity level $K\;(1\le K\le d)$
\For{$t=0,1,\dots,T-1$}
    \State Sample a minibatch (or a data point) $\xi_t\sim\mathcal D$
    \State Compute stochastic gradient $g_t=\nabla f(w_t;\xi_t)$
    \State $\displaystyle \tilde g_t\gets \topk{g_t}$ \Comment{keep only top $K$ entries}
    \State $w_{t+1}\gets w_t-\eta_t\tilde g_t$
\EndFor
\State \textbf{Return} $w_T$ (or the average $\bar w_T=\frac1T\sum_{t=0}^{T-1}w_t$)
\end{algorithmic}
\end{algorithm}

\section*{Dry Run Example}
Consider the scalar quadratic $f(w)=(w-3)^2$.
\[
\nabla f(w)=2(w-3).
\]
Set $w_0=0$, constant stepsize $\eta=0.1$, and $K=1$ (the only coordinate is kept).

\begin{center}
\begin{tabular}{c|c|c|c}
$t$ & $g_t= \nabla f(w_t)$ & $w_{t+1}=w_t-\eta\,\topk{g_t}$ & $f(w_{t+1})$\\\hline
0 & $2(0-3)=-6$ & $0-0.1(-6)=0.6$ & $(0.6-3)^2=5.76$\\
1 & $2(0.6-3)=-4.8$ & $0.6-0.1(-4.8)=1.08$ & $(1.08-3)^2=3.6864$\\
2 & $2(1.08-3)=-3.84$ & $1.08-0.1(-3.84)=1.464$ & $(1.464-3)^2=2.3650$
\end{tabular}
\end{center}
The objective value decreases at each iteration, illustrating the update (1) even though only the top‑$K$ component (the only one) is transmitted.

\newpage

\section*{Assumptions}
\begin{assumption}[Smoothness]\label{as:smooth}
$f$ is $L$‑smooth: for all $x,y\in\mathbb{R}^d$,
\[
f(y)\le f(x)+\langle\nabla f(x),y-x\rangle+\frac{L}{2}\|y-x\|^2 .
\]
\end{assumption}
\begin{assumption}[Unbiased Stochastic Gradient]\label{as:unbiased}
\[
\mathbb{E}_{\xi_t}[\,g_t\mid w_t\,]=\nabla f(w_t).
\]
\end{assumption}
\begin{assumption}[Bounded Variance]\label{as:variance}
There exists $\sigma^2>0$ such that
\[
\mathbb{E}_{\xi_t}\big[\|g_t-\nabla f(w_t)\|^2\mid w_t\big]\le\sigma^2 .
\]
\end{assumption}
\begin{assumption}[Top‑$K$ Compression Guarantee]\label{as:compress}
Define the compression factor
\[
\rho\;:=\;\frac{d-K}{K}\;\ge 0 .
\]
For any vector $v\in\mathbb{R}^d$ the top‑$K$ operator satisfies
\[
\|v-\topk{v}\|^2\;\le\;\rho\,\|v\|^2 . \tag{2}
\]
\end{assumption}
\begin{assumption}[Bounded Gradient Norm]\label{as:gradbound}
There exists $G>0$ such that $\|\nabla f(w)\|\le G$ for all $w$ visited by the algorithm.
\end{assumption}

\section*{Main Convergence Theorem}
\begin{theorem}\label{thm:main}
Assume \ref{as:smooth}–\ref{as:gradbound} hold and let the stepsize be constant
\[
\eta_t\equiv\eta\le\frac{1}{L(1+\rho)} .
\]
Then for any $T\ge1$,
\[
\frac{1}{T}\sum_{t=0}^{T-1}\mathbb{E}\big[\|\nabla f(w_t)\|^2\big]
\;\le\;
\frac{2\big(f(w_0)-f^\star\big)}{\eta T}
\;+\;
\eta L\sigma^2(1+\rho) .
\tag{3}
\]
Consequently, choosing $\eta=\Theta\big(\tfrac{1}{\sqrt{T}}\big)$ yields
\[
\frac{1}{T}\sum_{t=0}^{T-1}\mathbb{E}\big[\|\nabla f(w_t)\|^2\big]
=O\!\left(\frac{1}{\sqrt{T}}\right).
\]
\end{theorem}

\section*{Theoretical Properties}
\begin{itemize}
\item \textbf{Stability.} The condition $\eta\le1/[L(1+\rho)]$ guarantees that the extra error introduced by sparsification (captured by $\rho$) does not destabilize the recursion.
\item \textbf{Hyper‑parameter Sensitivity.} Larger $K$ (smaller $\rho$) permits a larger stepsize and improves the bound. In the limit $K=d$ we have $\rho=0$ and recover the standard SGD rate.
\item \textbf{Comparison to Baselines.} Compared with vanilla SGD, the additional term $\eta L\sigma^2\rho$ quantifies the slowdown due to communicating only $K$ coordinates. When $K\ll d$, $\rho\approx d/K$ and the degradation is linear in $d/K$.
\item \textbf{Special Cases.} If the stochastic gradients are exact ($\sigma=0$) the bound reduces to the deterministic rate $O(1/(\eta T))$, identical to gradient descent with a $1/(1+\rho)$‑scaled stepsize.
\end{itemize}

\section*{Discussion}
The theorem shows that Top‑K SGD retains the canonical $O(1/\sqrt{T})$ convergence of stochastic methods for non‑convex smooth objectives, at the cost of a multiplicative factor $1+\rho$ that reflects the information loss caused by sparsification. The result holds under minimal assumptions and directly parallels the classical SGD analysis, with the only modification being the compression inequality (2).

\newpage

\section*{Preliminaries (Definitions and Lemmas)}
\begin{lemma}[Smoothness Inequality]\label{lem:smooth}
If $f$ is $L$‑smooth then for any $x,y$,
\[
f(y)\le f(x)+\langle\nabla f(x),y-x\rangle+\frac{L}{2}\|y-x\|^2 .
\]
\end{lemma}
\begin{lemma}[Compression Error]\label{lem:compress}
For the top‑$K$ operator defined in \eqref{as:compress},
\[
\|v-\topk{v}\|^2\le\rho\|v\|^2,\qquad\rho=\frac{d-K}{K}.
\]
\end{lemma}
Both lemmas follow directly from the definitions and are used repeatedly in the proof.

\section*{Main Proof (Step‑by‑Step Derivation)}
We prove Theorem~\ref{thm:main}. Throughout, expectations are taken w.r.t. all randomness up to iteration $t$ unless otherwise noted.

\subsection*{Step 1 – One‑step descent}
From Lemma~\ref{lem:smooth} with $x=w_t$, $y=w_{t+1}=w_t-\eta\topk{g_t}$,
\[
f(w_{t+1})\le f(w_t)-\eta\langle\nabla f(w_t),\topk{g_t}\rangle
+\frac{L\eta^{2}}{2}\|\topk{g_t}\|^{2}. \tag{4}
\]

\subsection*{Step 2 – Decompose the inner product}
Add and subtract $g_t$ inside the inner product:
\[
\langle\nabla f(w_t),\topk{g_t}\rangle
=\langle\nabla f(w_t),g_t\rangle
+\langle\nabla f(w_t),\topk{g_t}-g_t\rangle . \tag{5}
\]

\subsection*{Step 3 – Take conditional expectation}
Condition on $w_t$ and use Assumption~\ref{as:unbiased}:
\[
\mathbb{E}\big[\langle\nabla f(w_t),g_t\rangle\mid w_t\big]
=\|\nabla f(w_t)\|^{2}. \tag{6}
\]

\subsection*{Step 4 – Bound the compression cross term}
Using Cauchy–Schwarz and Lemma~\ref{lem:compress},
\[
\big|\langle\nabla f(w_t),\topk{g_t}-g_t\rangle\big|
\le\|\nabla f(w_t)\|\,\|g_t-\topk{g_t}\|
\le\|\nabla f(w_t)\|\,\sqrt{\rho}\,\|g_t\| . \tag{7}
\]
Taking expectation and applying Young’s inequality $ab\le\frac{1}{2}(a^{2}+b^{2})$,
\[
\mathbb{E}\big[\langle\nabla f(w_t),\topk{g_t}-g_t\rangle\mid w_t\big]
\le\frac{1}{2}\|\nabla f(w_t)\|^{2}
+\frac{\rho}{2}\mathbb{E}\big[\|g_t\|^{2}\mid w_t\big]. \tag{8}
\]

\subsection*{Step 5 – Bound $\mathbb{E}\|g_t\|^{2}$}
Decompose $g_t=\nabla f(w_t)+(g_t-\nabla f(w_t))$ and use Assumption~\ref{as:variance}:
\[
\mathbb{E}\big[\|g_t\|^{2}\mid w_t\big]
= \|\nabla f(w_t)\|^{2}
+\mathbb{E}\big[\|g_t-\nabla f(w_t)\|^{2}\mid w_t\big]
\le \|\nabla f(w_t)\|^{2}+\sigma^{2}. \tag{9}
\]

\subsection*{Step 6 – Bound the squared norm of the compressed gradient}
From Lemma~\ref{lem:compress},
\[
\|\topk{g_t}\|^{2}
\le \|g_t\|^{2}. \tag{10}
\]
Taking expectation and using \eqref{9},
\[
\mathbb{E}\big[\|\topk{g_t}\|^{2}\mid w_t\big]
\le \|\nabla f(w_t)\|^{2}+\sigma^{2}. \tag{11}
\]

\subsection*{Step 7 – Plugging back into the descent inequality}
Insert \eqref{5}–\eqref{8} and \eqref{11} into \eqref{4}, then take conditional expectation:
\[
\begin{aligned}
\mathbb{E}\big[f(w_{t+1})\mid w_t\big]
&\le f(w_t)-\eta\|\nabla f(w_t)\|^{2}
+\eta\Big(\frac{1}{2}\|\nabla f(w_t)\|^{2}
+\frac{\rho}{2}\big(\|\nabla f(w_t)\|^{2}+\sigma^{2}\big)\Big)   \\
&\quad +\frac{L\eta^{2}}{2}\big(\|\nabla f(w_t)\|^{2}+\sigma^{2}\big).
\end{aligned}
\tag{12}
\]

\subsection*{Step 8 – Collect terms}
Grouping the coefficients of $\|\nabla f(w_t)\|^{2}$ and $\sigma^{2}$ yields
\[
\mathbb{E}\big[f(w_{t+1})\mid w_t\big]
\le f(w_t)
-\underbrace{\eta\Big(1-\tfrac12-\tfrac{\rho}{2}\Big)}_{=:c_1}\|\nabla f(w_t)\|^{2}
+\underbrace{\Big(\tfrac{\eta\rho}{2}+\tfrac{L\eta^{2}}{2}\Big)}_{=:c_2}\sigma^{2}
+\underbrace{\Big(\tfrac{L\eta^{2}}{2}\Big)}_{=:c_3}\|\nabla f(w_t)\|^{2}.
\tag{13}
\]

\subsection*{Step 9 – Use the stepsize restriction}
Choose $\eta\le\frac{1}{L(1+\rho)}$. Then
\[
c_3=\frac{L\eta^{2}}{2}\le\frac{\eta}{2(1+\rho)} .
\]
Consequently,
\[
c_1-c_3
=\eta\Big(\tfrac12-\tfrac{\rho}{2}\Big)-\frac{L\eta^{2}}{2}
\ge\eta\Big(\tfrac12-\tfrac{\rho}{2}-\tfrac{1}{2(1+\rho)}\Big)
\ge\frac{\eta}{2(1+\rho)} . \tag{14}
\]
Define
\[
\mu:=\frac{\eta}{2(1+\rho)}>0 .
\]

\subsection*{Step 10 – One‑step recursion}
Using \eqref{13} and the bound \eqref{14},
\[
\mathbb{E}\big[f(w_{t+1})\mid w_t\big]
\le f(w_t)-\mu\|\nabla f(w_t)\|^{2}+c_2\sigma^{2}.
\tag{15}
\]

\subsection*{Step 11 – Unconditional expectation and telescoping}
Take full expectation and sum \eqref{15} over $t=0,\dots,T-1$:
\[
\begin{aligned}
\mathbb{E}[f(w_T)]
&\le f(w_0)-\mu\sum_{t=0}^{T-1}\mathbb{E}\big[\|\nabla f(w_t)\|^{2}\big]
+T\,c_2\sigma^{2}.
\end{aligned}
\tag{16}
\]
Since $f(w_T)\ge f^\star$, rearrange:
\[
\mu\sum_{t=0}^{T-1}\mathbb{E}\big[\|\nabla f(w_t)\|^{2}\big]
\le f(w_0)-f^\star+T\,c_2\sigma^{2}.
\tag{17}
\]

\subsection*{Step 12 – Divide by $T$ and substitute constants}
Recall $c_2=\frac{\eta\rho}{2}+\frac{L\eta^{2}}{2}\le\frac{\eta}{2}(1+\rho)$ (using $\eta\le1/[L(1+\rho)]$). Hence,
\[
\frac{1}{T}\sum_{t=0}^{T-1}\mathbb{E}\big[\|\nabla f(w_t)\|^{2}\big]
\le \frac{f(w_0)-f^\star}{\mu T}
     +\frac{c_2}{\mu}\sigma^{2}.
\tag{18}
\]
Substituting $\mu=\frac{\eta}{2(1+\rho)}$ and $c_2\le\frac{\eta}{2}(1+\rho)$ gives
\[
\frac{1}{T}\sum_{t=0}^{T-1}\mathbb{E}\big[\|\nabla f(w_t)\|^{2}\big]
\le \frac{2\big(f(w_0)-f^\star\big)}{\eta T}
   +\eta L\sigma^{2}(1+\rho) .
\tag{19}
\]
Equation \eqref{19} is exactly the bound \eqref{3} claimed in Theorem~\ref{thm:main}. ∎

\section*{Rate Derivation (With Constants)}
From \eqref{19} choose $\eta = \frac{1}{\sqrt{T}}\frac{1}{\sqrt{L(1+\rho)}}$ (which respects the stepsize condition). Then
\[
\frac{1}{T}\sum_{t=0}^{T-1}\mathbb{E}\big[\|\nabla f(w_t)\|^{2}\big]
\le
\underbrace{\frac{2\big(f(w_0)-f^\star\big)\sqrt{L(1+\rho)}}{\sqrt{T}}}_{O(T^{-1/2})}
+
\underbrace{\frac{\sqrt{L(1+\rho)}}{\sqrt{T}}\,\sigma^{2}}_{O(T^{-1/2})}
=O\!\left(\frac{1}{\sqrt{T}}\right).
\]

\section*{Special Cases}
\begin{itemize}
\item \textbf{Full‑gradient (no compression).} $K=d\Rightarrow\rho=0$. The bound reduces to the classical SGD result $\frac{2(f(w_0)-f^\star)}{\eta T}+\eta L\sigma^{2}$.
\item \textbf{Deterministic setting.} $\sigma=0$ gives $\frac{2(f(w_0)-f^\star)}{\eta T}$, i.e. a $O(1/T)$ decay of the gradient norm, matching gradient descent with stepsize $1/(L(1+\rho))$.
\item \textbf{Extreme sparsity.} $K=1$ yields $\rho=(d-1)$; the stepsize must be reduced by a factor $1/(1+d-1)=1/d$, reflecting the severe information loss.
\end{itemize}

\end{document}